% !TeX spellcheck = de_DE
% !TeX program = lualatex

\DocumentMetadata{
  lang=de-DE,
  pdfstandard=ua-2,
  tagging = on,
  pdfversion=2.0,
  testphase={phase-III}
}

\documentclass[11pt,a4paper]{article}

\usepackage{fontspec}              
\usepackage[ngerman]{babel}
\usepackage[margin=12mm]{geometry}
\usepackage{graphicx}
\usepackage{xcolor}
\usepackage{hyperref}

\AssignTaggingSocketPlug{para/semantic/begin}{noop}
\AssignTaggingSocketPlug{para/semantic/end}{noop}

\setmainfont{Latin Modern Sans}
\setsansfont{Latin Modern Sans}
\renewcommand{\familydefault}{\sfdefault}

% Typografische Farbdefinitionen
\definecolor{cvblue}{HTML}{005580}
\definecolor{subtext}{HTML}{555555}

\tagpdfsetup{
  math/mathml/luamml/load=false,
  add-new-tag = Zeit/Span,
  add-new-tag = Ausbildungsort/Em,
  add-new-tag = Arbeitgeber/Em,
  add-new-tag = Ziel/Div,
  add-new-tag = Kategorie/Span,
  add-new-tag = Kenntnisse/Span,
  add-new-tag = Sprache/Span,
  add-new-tag = Niveau/Span,
  add-new-tag = KontaktLabel/Span,
  add-new-tag = KontaktWert/Span
}

\setcounter{secnumdepth}{0}

\newcommand{\tagtext}[2]{%
  \tagstructbegin{tag=#1}%
  \tagmcbegin{tag=#1}%
  #2%
  \tagmcend%
  \tagstructend%
}

% Beschreibungen enthalten Absätze und Listen. Kein Inline-MC um den Block!
% Die Rolle Ziel bleibt erhalten und wird auf den Blocktyp Div abgebildet.
\newcommand{\tagblock}[2]{%
  \tagstructbegin{tag=#1}%
  #2\par%
  \tagstructend%
}

\makeatletter
\renewcommand\section{\@startsection{section}{1}{\z@}%
                               {-0.2ex}%
                               {0.1ex}%
                               {\normalfont\Large\bfseries\color{cvblue}\scshape}}
\makeatother

\newcommand{\cvsection}[1]{%
  \section{#1}%
  \vspace*{-0.6ex}%
  \tagpdfparaOff
  \noindent{\color{cvblue}\rule{\linewidth}{0.8pt}}\par
  \vspace*{0.4ex}%
  \tagpdfparaOn
}

% Foto ist optional; ohne Foto nutzt die Kopfzeile die ganze Breite.
\newcommand{\cvheader}[4]{%
  \begin{minipage}[t]{0.25\textwidth}\vspace{0pt}\includegraphics[width=3.2cm,alt={Porträtfoto}]{\detokenize{/data/profilfoto.jpg}}\end{minipage}\hfill\begin{minipage}[t]{0.73\textwidth}\vspace{0pt}\raggedright
    {\Huge\bfseries\color{cvblue} #2}\par\vspace{0.4ex}
    {\large\bfseries\color{subtext} #3}\par\vspace{1ex}
    {\small #4}\par\end{minipage}
  \vspace{1ex}%
}

% Linke Spalte (Sidebar)
\newenvironment{kontakt}{\cvsection{Kontakt}}{}
\newlength{\kontaktlabelwidth}
\setlength{\kontaktlabelwidth}{1.85cm}
\newcommand{\kitem}[2]{%
  \noindent\hangindent=\kontaktlabelwidth\hangafter=1%
  \makebox[\kontaktlabelwidth][l]{\tagtext{KontaktLabel}{\textbf{#1}}}%
  \tagtext{KontaktWert}{#2}\par\vspace{0.4ex}%
}

\newenvironment{itkenntnisse}{\cvsection{IT-Kenntnisse}}{}
\newcommand{\ititem}[2]{%
  \noindent\tagtext{Kategorie}{\textbf{#1}}\newline
  {\small \tagtext{Kenntnisse}{#2}}\par\vspace{0.5ex}%
}

\newenvironment{sprachen}{\cvsection{Sprachen}}{}
\newcommand{\langitem}[2]{%
  \noindent\tagtext{Sprache}{\textbf{#1}} (\tagtext{Niveau}{#2})\par\vspace{0.4ex}%
}

% Rechte Spalte (Werdegang)
\newenvironment{ausbildung}{\cvsection{Ausbildung}}{}
\newcommand{\eduitem}[4]{%
  \noindent\textbf{#1}%
  \hfill {\small\color{subtext}\mbox{\tagtext{Zeit}{#2}}}\newline
  {\color{subtext}\tagtext{Ausbildungsort}{#3}}\par%
  \ifx&#4&\else
    {\small\tagblock{Ziel}{#4}}%
  \fi
  \vspace{1ex}%
}

\newenvironment{berufserfahrung}{\cvsection{Berufserfahrung}}{}
\newcommand{\jobitem}[4]{%
  \noindent\textbf{#1}%
  \hfill {\small\color{subtext}\mbox{\tagtext{Zeit}{#2}}}\newline
  {\color{subtext}\tagtext{Arbeitgeber}{#3}}\par%
  \ifx&#4&\else
    {\small\tagblock{Ziel}{#4}}%
  \fi
  \vspace{1ex}%
}

\hypersetup{
  colorlinks = true, 
  urlcolor = cvblue, 
  linkcolor = cvblue,
  pdftitle={Lebenslauf Alex Beispiel},
  pdfauthor={Alex Beispiel}
}

\begin{document}
\setlength{\topskip}{0pt}
\setlength{\parindent}{0pt}
\setlength{\parskip}{0pt}
\pagestyle{empty}

% Header
\cvheader{}
         {Alex Beispiel}
         {IT-Support und Systemadministration}
         {Fiktives Profil zur Demonstration. Alle Angaben sind erfunden.}

% Layout-Raster
\begin{minipage}[t]{0.33\textwidth}
\raggedright

  \begin{kontakt}
    \kitem{Adresse:}{\small Musterweg 10,\newline 12345 Musterstadt}
    \kitem{E-Mail:}{\href{mailto:alex@example.invalid}{\small alex@example.invalid}}
    \kitem{Telefon:}{\small +00 000 000000}
  \end{kontakt}%
  \vspace{0.3em}%
  %
  \begin{itkenntnisse}
    \ititem{Systeme}{Windows, Linux – Grundlagen}
    \ititem{Netzwerke}{TCP/IP, DNS, DHCP und VLANs}
    \ititem{Werkzeuge}{Python, Git und Virtualisierung}
  \end{itkenntnisse}%
  \vspace{0.3em}%
  %
  \begin{sprachen}
    \langitem{Deutsch}{Sehr gute Kenntnisse (Beispiel)}
    \langitem{Englisch}{Technische Dokumentation (Beispiel)}
  \end{sprachen}%
  %
\tagpdfsetup{sec/end=-10}%
\end{minipage}%
\hfill
\begin{minipage}[t]{0.64\textwidth}
\raggedright

  \begin{ausbildung}
    \eduitem{Lehrgang IT-Systemtechnik}
            {2023–2025}
            {Beispiel-Akademie, Musterstadt}
            {\begin{itemize}[left-margin=1.1em, label-sep=0.35em, label-width=0.75em, item-vspace=1pt, begin-vspace=2pt, end-vspace=2pt, begin-extra-vspace=0pt, end-extra-vspace=0pt, para-vspace=0pt]
    \item Windows- und Linux-Systeme in einer Testumgebung einrichten
    \item DNS- und Netzwerkfehler eingrenzen
    \item Konfigurationen und Lösungswege dokumentieren
\end{itemize}}
  \end{ausbildung}%
  \vspace{0.3em}%
  %
  \begin{berufserfahrung}
    \jobitem{Praktikum im IT-Support}
            {2024}
            {Example IT GmbH, Musterstadt}
            {\begin{itemize}[left-margin=1.1em, label-sep=0.35em, label-width=0.75em, item-vspace=1pt, begin-vspace=2pt, end-vspace=2pt, begin-extra-vspace=0pt, end-extra-vspace=0pt, para-vspace=0pt]
    \item Supportanfragen im Übungssystem bearbeiten
    \item Testgeräte einrichten und inventarisieren
\end{itemize}}
  \end{berufserfahrung}%
\cvsection{Projekte}
\eduitem{Virtuelles Testnetz}{2025}{Fiktives Lernprojekt}{\noindent Clients und Server in einem getrennten Netzwerk verbinden. Die Angaben dienen ausschließlich als Layoutbeispiel.\par}
  %
\tagpdfsetup{sec/end=-10}%
\end{minipage}

\end{document}
